\section{最近公共祖先}

\subsection{$\mathcal{O}(1)$ LCA}

使用片段：需上下文已有以 $1$ 为根的树邻接表 $G$（$1..n$）与前置模板 \texttt{SparseTable}。段内定义 \texttt{dfn}、\texttt{dep}、\texttt{fa}、\texttt{seq} 与查询函数 \texttt{lca(x,y)}。

预处理 $\mathcal{O}(n\log n)$，单次查询 $\mathcal{O}(1)$。

\begin{lstlisting}
vector<int> dfn(n+1),dep(n+1),fa(n+1),seq(1);
int dfnt = 0;
auto dfs = [&](this auto&& dfs,int u,int pa) -> void {
    dfn[u] = ++dfnt;
    dep[u] = dep[pa] + 1,fa[u] = pa;
    seq.push_back(u);
    for (int v : G[u]){
        if (v == pa) continue;
        dfs(v,u);
    }
};
dfs(1,1);
SparseTable st(seq, 1,[&](int x,int y) -> int {
    return dep[x] <= dep[y] ? x : y;
});
auto lca = [&](int x,int y) -> int {
    if (x == y) return x;
    x = dfn[x],y = dfn[y];
    if (x > y) swap(x,y);
    return fa[st(x+1,y)];
};
\end{lstlisting}

\subsection{树链剖分求 LCA}

使用片段：需上下文已有以 $1$ 为根的树邻接表 $G$（$1..n$）；段内定义 \texttt{dep}、\texttt{fa}、\texttt{son}、\texttt{top}、\texttt{dfn} 与查询函数 \texttt{lca(x,y)}。该版本为重链剖分。

预处理 $\mathcal{O}(n)$，单次查询时间复杂度为 $\mathcal{O}(\log n)$。

\begin{lstlisting}
vector<int> siz(n+1,1),dep(n+1),fa(n+1),son(n+1),top(n+1),dfn(n+1),seq(1);
auto dfs1 = [&](this auto&& dfs,int u,int pa)->void {
    fa[u] = pa;
    dep[u] = dep[pa] + 1;
    int mx = 0;
    for (int v : G[u]){
        if (v == pa) continue;
        dfs(v,u);
        siz[u] += siz[v];
        if (mx < siz[v]){
            son[u] = v;
            mx = siz[v];
        }
    }
};
int tick = 0;
auto dfs2 = [&](this auto&& dfs,int u,int tp)->void {
    dfn[u] = ++tick;
    seq.push_back(u);
    top[u] = tp;
    if (son[u]) {
        dfs(son[u],tp);
    }
    for (int v : G[u]){
        if (v == fa[u] || v == son[u]) continue;
        dfs(v,v);
    }
};
dfs1(1,1);
dfs2(1,1);
auto lca = [&](int x,int y) -> int {
    while (top[x] != top[y]) {
        if (dep[top[x]] > dep[top[y]]) swap(x,y);
        y = fa[top[y]];
    }
    return dep[x] > dep[y] ? y : x;
};
\end{lstlisting}

